Well-formedness rules out several Python runtime errors. A variable may be referenced only where it is
definitely assigned, and the fields and base class of a class may name only earlier type definitions, so
\pyinline{NameError} and \pyinline{UnboundLocalError} do not arise. An import requires its target to be a module
of the program, a from-import requires each imported name to be a member of the source module, and import
cycles are rejected, so \pyinline{ModuleNotFoundError} and \pyinline{ImportError} do not arise either. A module
attribute must resolve to a definitely assigned member and an object attribute must be a field of the class, which together rule out \pyinline{AttributeError},
and a constructor call or constructor pattern must match the arity of the class exactly, which rules out the
\pyinline{TypeError} raised by Python for excess arguments or sub-patterns. The semantics gives $\errName$,
$\errAttr$ and $\errType$ to the runtime errors ruled out by typing, and $\errExcluded$ to the statically excluded
behaviours (\secref{no-coercions}); type soundness, the conjecture that no run of a well-formed program aborts with
one of these kinds, is stated in \specref{sec:runtime-errors}.
